\paragraph*{Sampling strategy choice is subject-specific (panel validity).}
The choice of sampling strategies could itself drive the result. One concrete
hypothesis is that the cost-aware \texttt{pilotCost} sampler we add explains most
of the changes. We test this hypothesis with an experiment: we remove
\texttt{pilotCost} and the cost-aware bounds and re-run the comparison with only
the cost-unaware \texttt{random} and \texttt{balanced} samplers. The measured-cost
top group still changes in $59\%$ of \StudyEnergy, $52\%$ of \StudyWorkload, and
$32\%$ of \StudyLCBench\ comparisons (versus $50/46/42\%$ with the full set), the
count-recommended sampler is evicted in up to $22\%$, and Pareto fronts still
differ in $62$--$71\%$ of scenarios. The effect is thus a property of the cost
axis, not of the sampler we add, and the statistical verdicts are robust to
dropping it. Along the way the experiment teaches something more general: even
among cost-unaware samplers the cost-effective choice is subject- and
budget-specific, so practitioners should choose per subject rather than assume a
universal best. Consistently, the change also holds on the two replications and
\StudyLinux, which contain no cost-aware sampler and reuse the published
\StudyReplicationICSE\ and \StudyReplicationICPE\ strategy sets unchanged.

\paragraph*{Construct validity (what ``cost'' includes).}
We charge the measured cost of \emph{acquiring} the selected training
configurations, not the cost of \emph{selecting} them: constraint solving,
pairwise-distance computation, or active-learning retraining. This scope is
deliberate and conservative: prior sampling studies measure effort as sample
\emph{size} and thus ignore selection cost as well, so our cost axis is already a
strict improvement on the count baseline we compare against, and because most
strategies select a sample once and for all, any selection cost is amortized over
the campaign and small beside measurement.

\paragraph*{Internal validity (learners, seeds, grid).}
Prediction error, and thus which sampler wins RQ2 and RQ3, can depend on the
learner and on sampling noise (RQ1 does not, as it precedes any model fit). To
mitigate this we keep the replication studies' original ridge learner and re-run
the entire decision analysis under CART and the scikit-learn random forest
\cite{sklearn_api}, with gradient boosting on the high-dimensional Linux space,
which agrees with the random forest.
The strong learners cut error sharply (median MRE on \StudyReplicationICSE\
$0.26{\to}0.07$) yet do not restore the count proxy: count-versus-measured
top-choice agreement stays at $0.43$--$0.81$ and measured-cost Pareto fronts still
differ (Jaccard $0.30$--$0.57$). We average over repeated seeds (about $100$ runs
per cell for the replication archives and $30$ for the generated families) and
bootstrap every reported proportion ($B{=}1000$). Because sample sizes are
discrete, we judge decisions by Pareto changes across the size grid, not at a
single budget point.



\paragraph*{External validity.}
The systems span compilers, encoders, databases, archivers, energy, workloads,
machine-learning runtime, and Linux releases, all real configurable software
measured in prior studies rather than synthetic, though still a finite, curated
set, and we run many per-scope tests. We therefore report the families separately and rest
the claims on patterns that repeat across them rather than on any single test. A
system that recurs in two families is not a redundant replicate: it is measured on
a different non-functional property and context (x264, for instance, contributes
encoding time and output size over $17$ videos in \StudyReplicationICPE\ but
energy in \StudyEnergy), so agreement across families is evidence over genuinely
different measurement settings, not the same one repeated. Per-scope statistics,
$p$-values, and effect sizes are in the anonymized artifact~\cite{artifact}.
